\documentclass[12pt,a4paper]{article}
\usepackage[utf8]{inputenc}
\usepackage[T1]{fontenc}
\usepackage[french]{babel}
\usepackage[margin=2.5cm]{geometry}
\usepackage{hyperref}
\usepackage{graphicx}
\usepackage{xcolor}
\usepackage{listings}
\usepackage{booktabs}
\usepackage{tabularx}
\usepackage{titlesec}

% Setup colors for links and code
\definecolor{primary}{RGB}{180, 0, 0} % Palette rouge mentionnée dans le projet
\definecolor{codegray}{rgb}{0.5,0.5,0.5}
\definecolor{codepurple}{rgb}{0.58,0,0.82}
\definecolor{backcolour}{rgb}{0.95,0.95,0.92}

\hypersetup{
    colorlinks=true,
    linkcolor=primary,
    filecolor=magenta,      
    urlcolor=blue,
    pdftitle={Rapport de Projet - daibbar},
}

\lstdefinestyle{mystyle}{
    backgroundcolor=\color{backcolour},   
    commentstyle=\color{codegray},
    keywordstyle=\color{magenta},
    numberstyle=\tiny\color{codegray},
    stringstyle=\color{codepurple},
    basicstyle=\ttfamily\footnotesize,
    breakatwhitespace=false,         
    breaklines=true,                 
    captionpos=b,                    
    keepspaces=true,                 
    numbers=left,                    
    numbersep=5pt,                  
    showspaces=false,                
    showstringspaces=false,
    showtabs=false,                  
    tabsize=2
}
\lstset{style=mystyle}

\title{
    \vspace{2cm}
    \Huge \textbf{Projet daibbar} \\
    \Large Portfolio \& Blog Personnel \\
    \vspace{0.5cm}
    \large \textit{Learn, build, document, and share.}
    \vspace{2cm}
}
\author{
    \textbf{Réalisé par :} \\ DAIBBAR Mohamed \\ \textit{IID1 — ENSA Khouribga} \\[1cm]
    \textbf{Encadré par :} \\ Pr. Rabhi
}
\date{Année universitaire 2025 -- 2026}

\begin{document}

\maketitle
\thispagestyle{empty}
\newpage

\tableofcontents
\newpage

\section{Introduction et Présentation du Projet}

\subsection{Qu'est-ce que daibbar ?}
\textbf{daibbar} est un site web personnel doublé d'un moteur de blog léger. Son slogan officiel, \textit{« Learn, build, document, and share. »} (Apprendre, construire, documenter et partager), reflète l'essence même du projet. 

La page d'accueil suit un arc narratif clair et délibéré appelé \textbf{« Proof $\rightarrow$ Knowledge $\rightarrow$ Person »} :
\begin{itemize}
    \item \textbf{Hero} : Une section d'accroche avec le message \textit{Just Learn.}
    \item \textbf{Projets récents} : Une grille de présentation des réalisations (cartes avec des tuiles en dégradé, descriptions et technologies utilisées).
    \item \textbf{Derniers articles} : Un fil d'actualité sous forme de newsletter listant les dernières publications avec date, titre et aperçu.
    \item \textbf{About / Connect} : Une section biographique incluant des liens vers les réseaux professionnels (GitHub, LinkedIn, X, Email) agissant également comme pied de page.
\end{itemize}

\subsection{Objectifs du Projet}
Les principaux objectifs qui ont motivé la création de cette plateforme sont :
\begin{itemize}
    \item \textbf{Documenter} : Centraliser l'apprentissage et archiver les différents projets réalisés.
    \item \textbf{Partager} : Publier des articles techniques de manière accessible.
    \item \textbf{Présenter} : Mettre en valeur les compétences et les réalisations techniques dans un portfolio moderne.
    \item \textbf{Connecter} : Faciliter la prise de contact pour d'éventuelles collaborations.
\end{itemize}

\section{Architecture Technique}

Le projet a été développé en adoptant une stack technique moderne, robuste et performante, principalement axée sur l'écosystème JavaScript côté serveur.

\begin{table}[h]
\centering
\begin{tabularx}{\textwidth}{|l|X|}
\hline
\textbf{Couche} & \textbf{Technologie} \\ \hline
\textbf{Runtime} & Node.js 20 \\ \hline
\textbf{Serveur Web} & Express 5 \\ \hline
\textbf{Templates (Vues)} & EJS (Embedded JavaScript) \\ \hline
\textbf{Base de données} & PostgreSQL 16 \\ \hline
\textbf{Driver DB} & \texttt{pg} (avec gestion de pool de connexions) \\ \hline
\textbf{Conteneurisation} & Docker \& Docker Compose \\ \hline
\textbf{Frontend} & CSS Vanilla, Polices système, aucun framework \\ \hline
\textbf{Outils de développement} & \texttt{nodemon}, \texttt{dotenv} \\ \hline
\end{tabularx}
\caption{Stack technologique du projet daibbar}
\end{table}

\section{Structure du Projet}

L'application suit scrupuleusement le design pattern \textbf{MVC (Modèle-Vue-Contrôleur)} pour assurer une séparation propre des responsabilités.

\subsection{Arborescence}
Voici la structure principale du code source :

\begin{lstlisting}[language=bash, title=Structure des répertoires]
daibbar/
├── Dockerfile               # Image pour le conteneur de l'application
├── docker-compose.yml       # Orchestration App + PostgreSQL
├── package.json
├── .env.example
└── src/
    ├── server.js            # Point d'entrée de l'application Express
    ├── db/                  # Gestion de la base de données
    │   ├── pool.js          # Pool de connexion PostgreSQL
    │   └── populatedb.js    # Script d'initialisation de la DB (Seed)
    ├── routes/              # Routeurs Express
    │   ├── articles.js
    │   └── projects.js
    ├── controllers/         # Logique métier et handlers de requêtes
    │   ├── ArticlesController.js
    │   └── ProjectsController.js
    ├── model/               # Couche d'accès aux données
    │   ├── Article.js
    │   └── Project.js
    ├── views/               # Templates EJS (Rendu côté serveur)
    │   ├── landing.ejs      # Accueil complet
    │   ├── index.ejs        # Liste des articles
    │   └── About.ejs
    └── public/              # Fichiers statiques
        └── style.css        # Styles globaux
\end{lstlisting}

\subsection{Détails du Pattern MVC}
\begin{itemize}
    \item \textbf{Routes} : Elles définissent les endpoints HTTP de l'application (\texttt{/articles}, \texttt{/projects}).
    \item \textbf{Controllers} : Ils interceptent les requêtes, font appel aux modèles pour manipuler les données, et renvoient les vues correspondantes.
    \item \textbf{Model} : Encapsule les requêtes SQL via le pool \texttt{pg} pour interagir avec PostgreSQL de manière sécurisée.
    \item \textbf{Views} : Fichiers EJS générant le rendu HTML dynamique côté serveur.
\end{itemize}

\section{Fonctionnalités et Interface (UI/UX)}

\subsection{Design Minimaliste et Vanilla}
L'interface utilisateur a été conçue volontairement sans framework CSS volumineux. Le choix s'est porté sur du \textbf{CSS Vanilla} avec une approche esthétique très précise :
\begin{itemize}
    \item \textbf{Palette de couleurs} : Rouge, Blanc et Noir.
    \item \textbf{Typographie} : Style inspiré par Apple (gros titres avec un espacement négatif généreux pour les en-têtes, \textit{system font stack} et nombreux espaces blancs).
    \item \textbf{Navigation} : Barre de navigation collante (sticky) et translucide.
    \item Utilisation de variables personnalisées (CSS custom properties) pour gérer le thème global depuis le fichier \texttt{style.css}.
\end{itemize}

\subsection{Base de Données et Moteur d'Articles}
Actuellement, la base de données s'appuie sur une table principale pour stocker le contenu :
\begin{lstlisting}[language=SQL, title=Schéma de la table Article]
CREATE TABLE IF NOT EXISTS Article (
  id    SERIAL PRIMARY KEY,
  title TEXT,
  body  TEXT
);
\end{lstlisting}
Le script \texttt{src/db/populatedb.js} est fourni pour créer automatiquement cette table et y insérer des données de test (Seed) lors du premier déploiement.

\section{API et Routage}

L'application Express expose plusieurs routes pour servir les différentes pages :

\begin{table}[h]
\centering
\begin{tabularx}{\textwidth}{|l|l|l|X|}
\hline
\textbf{Méthode} & \textbf{Chemin} & \textbf{Handler / Contrôleur} & \textbf{Description} \\ \hline
GET & \texttt{/} & Rendu de \texttt{landing.ejs} & Affiche la page d'accueil complète. \\ \hline
GET & \texttt{/articles} & \texttt{ArticlesController} & Liste l'ensemble des articles issus de la base de données. \\ \hline
GET & \texttt{/About} & \textit{Stub} (En développement) & Page "À propos". \\ \hline
GET & \texttt{/Projects} & \textit{Stub} (En développement) & Page listant les projets. \\ \hline
ALL & \texttt{*} & 404 Handler & Retourne le message d'erreur "Route existe pas". \\ \hline
\end{tabularx}
\caption{Endpoints et Routage Express}
\end{table}

\section{Déploiement}

Le projet est conçu pour être déployé très facilement grâce à la conteneurisation.

\subsection{L'approche Docker (Recommandée)}
L'utilisation de \texttt{docker-compose} permet de démarrer l'application Node.js (exposée sur le port 3000) ainsi que la base de données PostgreSQL (exposée sur le port 5432) de manière conjointe. 

\begin{lstlisting}[language=bash, title=Lancement via Docker]
# Cloner le dépôt et configurer les variables d'environnement
cp .env.example .env

# Lancer les conteneurs en tâche de fond
docker compose up

# Initialiser la base de données (au premier lancement uniquement)
docker compose exec app node src/db/populatedb.js
\end{lstlisting}

\subsection{Variables d'Environnement}
Un fichier \texttt{.env} à la racine du projet est nécessaire. Il inclut les valeurs de configuration suivantes :
\begin{itemize}
    \item \texttt{DATABASE\_URL} : Chaîne de connexion complète PostgreSQL (lue par \texttt{pool.js}).
    \item \texttt{DB\_USER}, \texttt{DB\_PASS}, \texttt{DB\_NAME} : Paramètres utilisés pour provisionner l'instance Postgres au sein du conteneur.
\end{itemize}

\section{Roadmap et Perspectives}

Le projet s'inscrit dans une démarche d'amélioration continue avec plusieurs évolutions prévues.

\subsection{Version 2 (V2)}
\begin{itemize}
    \item Connexion complète de la page d'accueil aux données réelles de la base de données (pour le fil d'articles).
    \item Implémentation complète des pages \texttt{/About} et \texttt{/Projects} avec de vraies réalisations, en remplaçant les placeholders actuels.
    \item Création d'une interface d'administration CRUD (Create, Read, Update, Delete) complète pour les articles.
    \item Mise en place d'une authentification sécurisée limitant l'accès aux actions d'administration.
\end{itemize}

\subsection{Version 3 (V3)}
\begin{itemize}
    \item Ajout d'un système de mots-clés (tags) et d'un filtrage des articles par catégories.
    \item Intégration d'un mode Sombre/Clair dynamique (Dark Mode) en sauvegardant les préférences de l'utilisateur.
\end{itemize}

\section{Conclusion}

Le projet \textbf{daibbar} illustre la conception rigoureuse d'une plateforme web sur-mesure combinant un portfolio personnel et un moteur de blog léger. En s'appuyant sur des technologies back-end fiables (Node.js, Express, PostgreSQL) couplées à une approche front-end \textit{Vanilla} minimaliste assumée, le site bénéficie d'une architecture MVC claire, facilement maintenable et déployable (via Docker). Ce rapport démontre que les bases solides actuelles laissent place à de belles perspectives d'évolution (V2 et V3) pour enrichir continuellement ses fonctionnalités.

\end{document}